% ============================================================================
% LANGENTWURF LE-016: Textverarbeitung Basics
% Profil: S (Senioren 65+) | Tier: 2 (Standard)
% ============================================================================

\documentclass[11pt,a4paper]{article}
\usepackage[utf8]{inputenc}
\usepackage[T1]{fontenc}
\usepackage[ngerman]{babel}
\usepackage{geometry}
\usepackage{fancyhdr}
\usepackage{lastpage}
\usepackage{xcolor}
\usepackage{tcolorbox}
\usepackage{enumitem}
\usepackage{longtable}
\usepackage{hyperref}
\usepackage{amssymb}

\geometry{left=2.5cm, right=2.5cm, top=2.5cm, bottom=2.5cm}
\definecolor{fach}{HTML}{b35c00}
\definecolor{gray}{HTML}{666666}
\definecolor{successgreen}{HTML}{2a7a4b}

\tcbuselibrary{skins,breakable}
\newtcolorbox{scorebox}{ colback=white, colframe=fach, fonttitle=\bfseries, title=Didaktik-Score, breakable }
\newtcolorbox{tippbox}[1][]{ colback=successgreen!5, colframe=successgreen, fonttitle=\bfseries, title=#1, breakable }

\pagestyle{fancy}
\fancyhf{}
\fancyhead[L]{\small\textcolor{gray}{Digitale Grundlagen -- 016 Textverarbeitung}}
\fancyhead[R]{\small\textcolor{gray}{LE Tier 2 / Profil S}}
\fancyfoot[C]{\small\thepage{} / \pageref{LastPage}}
\renewcommand{\headrulewidth}{0.4pt}

\begin{document}

\begin{titlepage}
    \centering
    \vspace*{2cm}
    {\Large\textcolor{gray}{Digitale Grundlagen -- Senioren 65+}}\\[2cm]
    {\Huge\textcolor{fach}{\textbf{Langentwurf}}}\\[0.5cm]
    {\large Tier 2 | Profil S}\\[2cm]
    {\LARGE\textbf{016 -- Textverarbeitung Basics}}\\[0.5cm]
    {\large Modul D: Produktivität}\\[2cm]
    \begin{tabular}{rl}
        \textbf{Zeitrahmen:} & 60--75 Minuten \\
        \textbf{Vorkenntnisse:} & Tippen auf Bildschirmtastatur \\
    \end{tabular}
    \vfill
    {\normalsize Stand: \today}
\end{titlepage}

\tableofcontents
\newpage

\section{Bedingungsanalyse}

\subsection{Ausgangssituation}
\begin{itemize}
    \item Notizen-App ist oft unbekannt oder ungenutzt
    \item Zettelwirtschaft statt digitaler Notizen
    \item Diktierfunktion ist unbekannt
    \item Unterschied: Notizen vs. Pages vs. Word unklar
    \item Texte werden nicht gespeichert/gefunden
\end{itemize}

\subsection{Typische Probleme}
\begin{itemize}
    \item ``Wo schreibe ich das auf?''
    \item ``Wie finde ich meine Notiz wieder?''
    \item ``Kann ich auch diktieren?''
    \item ``Wie drucke ich das aus?''
\end{itemize}

\section{Sachanalyse}

\subsection{Kernkonzepte}
\begin{enumerate}
    \item \textbf{Notizen-App:} Einfache Texte schreiben und organisieren
    \item \textbf{Ordner:} Notizen in Kategorien sortieren
    \item \textbf{Diktierfunktion:} Sprechen statt Tippen
    \item \textbf{Teilen:} Notiz per Mail oder WhatsApp verschicken
    \item \textbf{iCloud-Sync:} Notizen auf allen Geräten verfügbar
\end{enumerate}

\subsection{Alltagsanalogien}
\begin{tabular}{|l|p{8cm}|}
\hline
\textbf{Begriff} & \textbf{Analogie} \\
\hline
Notizen-App & Digitaler Notizblock, der nie voll wird \\
Ordner & Verschiedene Notizblöcke für verschiedene Themen \\
Diktierfunktion & Diktiergerät mit automatischer Schreibkraft \\
Sync & Kopie erscheint automatisch auf allen Geräten \\
\hline
\end{tabular}

\section{Didaktische Analyse}

\subsection{Stellung in der Reihe}
Erste Produktivitäts-Einheit. Grundlage für Organisation und Dokumentation.

\subsection{Didaktische Reduktion}
\textbf{Ausgeklammert:} Pages/Word, komplexe Formatierung, Tabellen, Checklisten-Funktion.

\textbf{Fokus:} Notizen-App für einfache Texte und Listen.

\section{Lernziele}

\begin{tabular}{|c|p{7cm}|c|c|}
\hline
\textbf{Nr.} & \textbf{Die Teilnehmer können...} & \textbf{Stufe} & \textbf{Niveau} \\
\hline
FZ1 & die Notizen-App öffnen und eine neue Notiz erstellen & Anwenden & Basis \\
FZ2 & Text eingeben (tippen oder diktieren) & Anwenden & Basis \\
FZ3 & eine Notiz wiederfinden (Suche) & Anwenden & Basis \\
FZ4 & Ordner erstellen und Notizen sortieren & Anwenden & Vertiefung \\
FZ5 & eine Notiz teilen (Mail, WhatsApp) & Anwenden & Vertiefung \\
\hline
\end{tabular}

\section{Methodische Analyse}

\subsection{Diktierfunktion nutzen}
\begin{tippbox}[Sprechen statt Tippen]
\begin{enumerate}
    \item Mikrofon-Symbol auf der Tastatur antippen
    \item Deutlich und langsam sprechen
    \item Satzzeichen mitsprechen: ``Punkt'', ``Komma'', ``Fragezeichen''
    \item ``Neue Zeile'' für Absatz
    \item Erneut antippen zum Beenden
\end{enumerate}
\end{tippbox}

\section{Verlaufsplanung}

\begin{longtable}{|p{1cm}|p{2cm}|p{5cm}|p{3cm}|p{2cm}|}
\hline
\textbf{Zeit} & \textbf{Phase} & \textbf{Inhalt} & \textbf{TN-Aktivität} & \textbf{Material} \\
\hline
0--10 & Einstieg & ``Wo schreiben Sie sich etwas auf?'' Zettel vs. digital & Austauschen & -- \\
\hline
10--25 & Notizen-App & App öffnen, Übersicht, neue Notiz erstellen & Mitmachen & S.1 \\
\hline
25--40 & Text eingeben & Tippen üben, Diktierfunktion vorstellen & Selbst üben & S.1 \\
\hline
40--55 & Organisieren & Suche, Ordner erstellen, Notizen verschieben & Üben & S.2 \\
\hline
55--70 & Teilen & Notiz per Mail/WhatsApp senden & Ausprobieren & S.2-3 \\
\hline
\end{longtable}

\section{Lösungen}

\textbf{Neue Notiz erstellen:}
\begin{itemize}
    \item Notizen-App öffnen (gelbes Symbol mit Stift)
    \item Unten rechts: Quadrat mit Stift (neue Notiz)
    \item Sofort losschreiben -- erste Zeile wird Titel
\end{itemize}

\textbf{Notiz finden:}
\begin{itemize}
    \item Oben: Suchfeld antippen
    \item Stichwort eingeben
    \item Alle Notizen mit diesem Wort werden angezeigt
\end{itemize}

\textbf{Ordner erstellen:}
\begin{itemize}
    \item In der Ordnerübersicht (ganz links)
    \item Unten links: ``Neuer Ordner''
    \item Namen eingeben (z.B. ``Einkaufslisten'', ``Arzt'')
\end{itemize}

\textbf{Praktische Notiz-Ideen:}
\begin{itemize}
    \item Einkaufsliste
    \item Telefonnummern für Notfälle
    \item Medikamenten-Plan
    \item Passwort-Hinweise (nicht die Passwörter selbst!)
    \item Geburtstage
\end{itemize}

\section{Didaktik-Score}

\begin{scorebox}
\begin{tabular}{|c|l|c|c|p{3.5cm}|}
\hline
\# & Dimension & Gew. & Score & Notiz \\
\hline
1 & Differenzierung & 15\% & 8/10 & Basis/Vertiefung klar \\
2 & Taxonomie & 10\% & 9/10 & Anwenden dominant \\
3 & Alltagsrelevanz & 10\% & 9/10 & Praktisch nützlich \\
4 & Aktivierung & 10\% & 10/10 & Eigene Notizen \\
5 & Scaffolding & 10\% & 9/10 & Schritt-für-Schritt \\
6 & Feedback & 10\% & 8/10 & Checkliste vorhanden \\
7 & Sprachliche Klarheit & 10\% & 9/10 & Einfache Sprache \\
8 & Motivation & 10\% & 9/10 & Ersetzt Zettelwirtschaft \\
9 & Barrierefreiheit & 10\% & 10/10 & Diktierfunktion! \\
10 & Praktikabilität & 5\% & 9/10 & 70 Min gut \\
\hline
\multicolumn{3}{|r|}{\textbf{GESAMT:}} & \textbf{9.0/10} & \textcolor{successgreen}{\textbf{FREIGABE}} \\
\hline
\end{tabular}
\end{scorebox}

\end{document}
